\documentclass[12pt,a4paper]{article}

% Balíčky pro češtinu a kódování
\usepackage[utf8]{inputenc}
\usepackage[T1]{fontenc}
\usepackage[czech]{babel}
\usepackage[margin=2.0cm]{geometry}
\usepackage{xcolor}
\usepackage{hyperref}
\usepackage{tabularx}
\usepackage{booktabs}
\usepackage{listings}

% Nastavení pro výpisy PL/SQL kódu
\definecolor{codegreen}{rgb}{0,0.6,0}
\definecolor{codegray}{rgb}{0.5,0.5,0.5}
\definecolor{codepurple}{rgb}{0.58,0,0.82}
\definecolor{backcolour}{rgb}{0.95,0.95,0.92}

\lstdefinestyle{mystyle}{
    backgroundcolor=\color{backcolour},   
    commentstyle=\color{codegreen},
    keywordstyle=\color{blue}\bfseries,
    numberstyle=\tiny\color{codegray},
    stringstyle=\color{codepurple},
    basicstyle=\ttfamily\footnotesize,
    breakatwhitespace=false,         
    breaklines=true,                 
    captionpos=b,                    
    keepspaces=true,                 
    numbers=left,                    
    numbersep=5pt,                  
    showspaces=false,                
    showstringspaces=false,
    showtabs=false,                  
    tabsize=2,
    language=SQL
}
\lstset{style=mystyle, inputencoding=utf8}

\title{\textbf{Ultimátní PL/SQL Referenční Příručka}}
\author{Komplexní přehled syntaxe a řešení chyb}
\date{\today}

\begin{document}

\maketitle
\tableofcontents
\newpage

\section{Systémový katalog a obecná pravidla}
\begin{itemize}
    \item \textbf{Katalog (USER\_TABLES):} Názvy jsou v databázi zásadně \textbf{VELKÝMI PÍSMENY}. Při dynamickém hledání vždy používej \texttt{UPPER(p\_name)}.
    \item \textbf{Apostrofy vs. Uvozovky:}
    \begin{itemize}
        \item \texttt{' '} (Jednoduché): Pro textové řetězce (např. \texttt{'REPORT'}).
        \item \texttt{" "} (Dvojité): Pro vynucení názvu objektu (v kódu se jim vyhýbej).
    \end{itemize}
    \item \textbf{Správa místa:} Pokud dojde kvóta (ORA-01652), použij \texttt{PURGE RECYCLEBIN;} nebo maž tabulky přes \texttt{DROP TABLE nazev PURGE;}.
    \item \textbf{Operátor LIKE:} Vyžaduje zástupný znak \texttt{\%} (např. \texttt{LIKE 'J\_REPORT\%'}).
\end{itemize}

\section{SQL Specifika a Analýza dotazů}
\begin{itemize}
    \item \textbf{Pořadí zpracování (Problém Alias):} Klauzule \texttt{HAVING} se vyhodnocuje \textbf{před} klauzulí \texttt{SELECT}. Proto v \texttt{HAVING} nelze použít alias (např. \texttt{article\_count}). Musíš tam zopakovat celou agregační funkci.
    \item \textbf{CTAS (Create Table As Select):} Každý vypočítaný sloupec \textbf{musí} mít alias (\texttt{AS jmeno}), jinak vytvoření tabulky selže na chybě \texttt{ORA-00998}.
    \item \textbf{Filtrování u LEFT JOIN:} Podmínka pro levou tabulku ve \texttt{WHERE} změní spojení na \texttt{INNER JOIN}. Pro zachování prázdných řádků musí být podmínka v klauzuli \texttt{ON}.
    \item \textbf{Agregace a SYSTIMESTAMP:} Funkce jako \texttt{SYSTIMESTAMP} nebo \texttt{SYSDATE} se nedávají do \texttt{GROUP BY}.
    \item \textbf{SQL\%ROWCOUNT:} Systémový atribut, který vrací počet řádků ovlivněných posledním \texttt{UPDATE/DELETE/INSERT}. Kontroluj hned po příkazu přes \texttt{IF SQL\%ROWCOUNT = 0 THEN...}.
\end{itemize}

\section{Dynamické SQL a Proměnné}
\begin{itemize}
    \item \textbf{Statické SQL:} Uvnitř \texttt{BEGIN...END} musí mít každý \texttt{SELECT} klauzuli \texttt{INTO}.
    \item \textbf{Vázané proměnné (USING):} Lze použít \textbf{pouze} v DML příkazech pro hodnoty.
    \item \textbf{Zákaz v DDL:} V příkazech \texttt{CREATE}, \texttt{DROP} nebo \texttt{ALTER} nelze použít \texttt{USING}. Hodnoty se musí lepit přímo do řetězce přes \texttt{||}.
    \item \textbf{Deklarace typu \%TYPE / \%ROWTYPE:}
    \begin{itemize}
        \item \texttt{v\_jmeno z\_author.name\%TYPE;} (okopíruje typ jednoho sloupce).
        \item \texttt{v\_radek z\_article\%ROWTYPE;} (vytvoří proměnnou se strukturou celého řádku).
    \end{itemize}
\end{itemize}

\section{Datové modelování (Oracle Academy)}
\begin{itemize}
    \item \textbf{Vztah 1:1 (Optional to Mandatory):} Cizí klíč (FK) patří \textbf{vždy na povinnou stranu} (Tabulka B).
    \item \textbf{Složený PK:} \texttt{ALTER TABLE x ADD CONSTRAINT pk PRIMARY KEY (col1, col2);}.
\end{itemize}

\newpage
\section{Výjimky (Exceptions)}

\begin{table}[h]
\centering
\begin{tabularx}{\textwidth}{|l|l|X|}
\hline
\textbf{Jméno výjimky} & \textbf{Číslo chyby} & \textbf{Popis} \\ \hline
\texttt{NO\_DATA\_FOUND} & ORA-01403 & \texttt{SELECT INTO} nevrátil nic. (\texttt{COUNT} ji nevyvolá). \\ \hline
\texttt{TOO\_MANY\_ROWS} & ORA-01422 & \texttt{SELECT INTO} vrátil více řádků. \\ \hline
\texttt{DUP\_VAL\_ON\_INDEX} & ORA-00001 & Porušení unikátního klíče. \\ \hline
\texttt{VALUE\_ERROR} & ORA-06502 & Chyba velikosti nebo typu (např. dlouhý string). \\ \hline
\end{tabularx}
\end{table}

\noindent
\begin{minipage}{\textwidth}
\subsection*{Uživatelské výjimky}
\begin{lstlisting}[language=SQL]
DECLARE
    e_moje_chyba EXCEPTION;
BEGIN
    IF neco_je_spatne THEN RAISE e_moje_chyba; END IF;
EXCEPTION
    WHEN e_moje_chyba THEN dbms_output.put_line('Chyba logiky!');
    WHEN OTHERS THEN dbms_output.put_line(SQLERRM);
END;
\end{lstlisting}
\end{minipage}

\vspace{1em}
\noindent
\begin{minipage}{\textwidth}
\subsection*{Výjimka s vlastním kódem}
Rozsah čísel musí být \textbf{-20000 až -20999}.
\begin{lstlisting}[language=SQL]
BEGIN
    raise_application_error(-20001, 'Vlastni text chyby pro uzivatele');
END;
\end{lstlisting}
\end{minipage}

\newpage
\section{Šablony objektů}

\noindent
\begin{minipage}{\textwidth}
\subsection{1. PROCEDURE (Procedura)}
\begin{lstlisting}[language=SQL]
CREATE OR REPLACE PROCEDURE P_priklad(p_id INT, p_mode VARCHAR2) IS
    v_sql VARCHAR2(2000);
    v_exists INT;
BEGIN
    SELECT count(*) INTO v_exists FROM user_tables WHERE table_name = 'TEST';
    
    IF p_mode = 'REPLACE' AND v_exists > 0 THEN
        EXECUTE IMMEDIATE 'DROP TABLE TEST';
    END IF;
    
    v_sql := 'CREATE TABLE TEST AS SELECT * FROM z_article WHERE aid = ' || p_id;
    EXECUTE IMMEDIATE v_sql;
END P_priklad;
\end{lstlisting}
\end{minipage}

\vspace{2em}
\noindent
\begin{minipage}{\textwidth}
\subsection{2. FUNCTION (Funkce)}
\begin{lstlisting}[language=SQL]
CREATE OR REPLACE FUNCTION F_pocet(p_year INT) RETURN INT IS
    v_res INT;
BEGIN
    SELECT count(*) INTO v_res FROM z_article WHERE year = p_year;
    RETURN v_res;
EXCEPTION
    WHEN OTHERS THEN RETURN -1;
END F_pocet;
\end{lstlisting}
\end{minipage}

\vspace{2em}
\noindent
\begin{minipage}{\textwidth}
\subsection{3. TRIGGER (Kombinovaný s predikáty)}
\begin{lstlisting}[language=SQL]
CREATE OR REPLACE TRIGGER T_hlidac
    BEFORE INSERT OR UPDATE OR DELETE ON z_article_author FOR EACH ROW
BEGIN
    IF INSERTING OR UPDATING THEN
        IF :NEW.rid IS NULL THEN raise_application_error(-20005, 'RID chybi'); END IF;
        :NEW.last_change := SYSTIMESTAMP; -- Prirazeni do :NEW
    END IF;
    
    IF DELETING THEN
        dbms_output.put_line('Mazu autora: ' || :OLD.rid);
    END IF;
END T_hlidac;
\end{lstlisting}
\end{minipage}

\vspace{2em}
\noindent
\begin{minipage}{\textwidth}
\subsection{4. FOR LOOP (Implicitní kurzor)}
\begin{lstlisting}[language=SQL]
BEGIN
    FOR r IN (SELECT DISTINCT year FROM z_article ORDER BY year) LOOP
        dbms_output.put_line('Zpracovavam rok: ' || r.year);
        -- Dynamicke SQL uvnitr cyklu
        EXECUTE IMMEDIATE 'CREATE TABLE R_' || r.year || ' AS SELECT...';
    END LOOP;
END;
\end{lstlisting}
\end{minipage}

\end{document}
